\section{Introducción}

La generación de residuos agrícolas constituye un desafío ambiental debido a la cantidad de biomasa que puede ser desaprovechada o manejada de manera inadecuada. En este contexto, la cáscara de papa representa un residuo de interés debido a su disponibilidad y a la posibilidad de aprovechar su contenido de materia orgánica mediante procesos de transformación termoquímica.

El presente proyecto de investigación plantea el aprovechamiento y valorización de cáscara de papa de producción boyacense mediante la producción de biocarbón (biochar) a través del proceso de pirólisis. Se busca estudiar las condiciones de producción del material obtenido y explorar su potencial de aplicación agronómica y ambiental.

La investigación parte de la necesidad de transformar un residuo agrícola en un producto con potencial utilidad, contribuyendo al aprovechamiento de recursos y a la búsqueda de alternativas para el manejo de residuos orgánicos. Durante el desarrollo del proyecto se pretende recopilar y analizar información relacionada con las características de la materia prima, las condiciones del proceso de pirólisis y las propiedades del biochar obtenido.

Esta introducción corresponde a la formulación inicial de la idea de investigación y podrá ser ampliada y modificada conforme avance el proyecto y se desarrollen las diferentes etapas metodológicas.
